\documentclass[10pt,a4paper]{amsart}
\usepackage[margin=19mm]{geometry}
\usepackage{amsmath,amssymb,mathtools}
\usepackage[scr=rsfs,cal=euler]{mathalfa}
\usepackage{xfrac}
\usepackage[unicode,hidelinks,bookmarks=false]{hyperref}
\newcommand{\ie}{i.e.\ }
\newcommand{\cf}{cf.\ }
\newcommand{\resp}{respectively\ }
\newcommand{\Subsection}{\subsection}
\providecommand{\nobreakdash}{\nobreak\hbox{-}}
\setlength{\emergencystretch}{2em}
\multlinegap0pt
\usepackage{amssymb}
\usepackage{xcolor}
\usepackage{algorithm}\usepackage{algpseudocode}
\usepackage{float}
\newtheorem{lemma}{Lemma}
\newcommand{\DG}{\mathrm{DG}}
\newcommand{\bfU}{{\boldsymbol U}}
\newcommand\numberthis{\addtocounter{equation}{1}\tag{\theequation}}
\title{Numerical Analysis of a Coupled 3D-1D Transport Problem}
\author{Alyssa Taylor-LaPole, Uzochi Gideon, Beatrice Riviere, Duygu Vargun}
\date{Source: arXiv:2603.19395v1 (CC BY 4.0)}
\begin{document}
\maketitle

\noindent\emph{Source and licence.} Numerical Analysis of a Coupled 3D-1D Transport Problem. Version arXiv:2603.19395v1 (CC BY 4.0), distributed by arXiv under the Creative Commons Attribution 4.0 International licence, http://creativecommons.org/licenses/by/4.0/, as recorded in the arXiv licence metadata for this exact version and shown on the abstract page https://arxiv.org/abs/2603.19395v1. Full licence notice: \texttt{licenses/arxiv-2603.19395-CC-BY-4.0.txt}.\par\smallskip

\noindent\textit{Editorial scope.} Counted unit: the article's lemma eq:stability (source0.tex lines 420--438). The article's complete original proof of it is delivered with it (source0.tex lines 440--562, one \texttt{proof} environment of 123 lines). No mathematics is written, repaired or re-derived: the statement and the proof are the article's own lines, in the article's own notation, and no retained line is substituted or re-flowed.  Retained uncounted as the source-local context the proof needs: lines 276--279; lines 362--379.  Nothing was omitted from any retained span, so the package carries no omission record.  No retained line contains a citation, so the wrapper carries no bibliography and no bibliography entry.  No retained line contains a figure environment, an image-inclusion command or drawing code, so no image is delivered and the package declares no asset dependency.  Editorial formatting only, in addition: an AMS \texttt{amsart} preamble that restates the article's own declarations and macros used by the retained lines, namely the environments \texttt{lemma} on separate counters, together with the standard packages those lines need.  The retained environments print in this document as Lemma 1 in order of appearance, whereas the article prints the same items with the numbering of its own full document.  The complete native source of the article as distributed in the arXiv e-print for this exact version is bundled unchanged as \texttt{sources/196767/source0.tex}.
    \begin{align}
    \mbox{Case 1:}&\quad \nabla \cdot \bfU = 0, \quad \eta_\bfU = 0,\\
    \mbox{Case 2:}&\quad \Vert \bfU \Vert_{L^\infty(\Omega)} \leq \frac{k_0}{2 C_0}, \quad \eta_\bfU = 1.\label{eq:condU}
    \end{align}

\begin{align}
    \frac{1}{\tau}\int_\Omega (c_h^{n}-c_h^{n-1}) v_h
      +  & \int_\Omega  \kappa \nabla c_h^{n} \cdot \nabla v_h  
       - \int_\Omega \bfU c_h^{n} \cdot \nabla v_h
   \nonumber\\
      +&\int_\Lambda\gamma \vert \partial\mathcal{D}\vert \left(\overline{c_h^{n}}-\hat{c}_h^{n}\right)\overline{v_h}
        = \int_\Omega f(\cdot,t^{n}) v_h, \quad \forall v_h \in V_h^\Omega,\label{eq:FEMscheme}\\
    \frac{1}{\tau} \int_\Lambda |\mathcal{D}|(\hat c_h^{n}- \hat c_h^{n-1}) \hat v_h 
    + & a_\Lambda(\hat c_h^{n}, \hat v_h)
    + b_\Lambda(\hat c_h^{n}, \hat v_h)
    +\int_\Lambda \gamma|\partial\mathcal{D}|\left(\hat c_h^{n}-\overline{c_h^{n}}\right)   \hat v_h
    \nonumber\\
   =  &  \int_\Lambda \hat f(\cdot, t^{n}) \hat v_h
   + |\mathcal{D}(0)|\hat U \, c^{\mathrm{in}}(t^n) \hat v_h(0),
   \quad \forall \hat v_h \in V_h^\Lambda,\label{eq:DGscheme}
   % \\- \theta_0 v_1(s_0) + \theta_N v_1|_{s_N}\numberthis.
    %\label{eq:1DWFdG}
\end{align}

\begin{lemma}
    For each $n\geq 1$, there exists a unique solution $(c_h^{n}, \hat c_h^{n})\in V_h^\Omega\times V_h^\Lambda$ satisfying \eqref{eq:FEMscheme}-\eqref{eq:DGscheme}.   In addition, %if $c^0\in H^s(\Omega)$ with $s>3/2$, $\hat c^0 \in H^r(\Omega$ with $r>1/2$ and 
    if $\tau \leq 1/2$, there is a constant $M>0$ independent of $h$ and $\tau$ such that
\begin{align}
\forall m\geq 1, \quad 
    \Vert c_h^m\Vert_{L^2(\Omega)}^2 
    & +d_0 \Vert \hat c_h^n\Vert_{L^2(\Lambda)}^2 
      +  \tau k_0 \sum_{n=1}^m \Vert \nabla c_h^n \Vert_{L^2(\Omega)}^2
      + \tau C_a \sum_{n=1}^m \vert \hat c_h^n \vert_{\DG}^2 \nonumber\\
        \leq  
        & \, M \tau \sum_{n=1}^m \Vert f(\cdot,t^n)\Vert_{L^2(\Omega)}^2
   + M \tau \sum_{n=1}^m \Vert \hat f(\cdot,t^n) \Vert_{L^2(\Lambda)}^2 
   \nonumber\\
&   + M |\mathcal{D}(0)|\hat U \, \tau \sum_{n=1}^{m}(c^{\mathrm{in}}(t^n))^2
+ C \Vert c^{0}\Vert_{H^s(\Omega)}^2 
+ M \Vert \hat c^{0}\Vert_{H^r(\Lambda)}^2.
\label{eq:stability}
\end{align}
\end{lemma}

\begin{proof}
Assume that $n\geq 1$.
Since the problem is linear and finite-dimensional, it suffices to show uniqueness of the solutions.  Assume that $(w,\hat w) = (c_h^{n,1}-c_h^{n,2}, \hat c_h^{n,1}-\hat c_h^{n,2})$, is the pair of differences between two solution pairs. It suffices to show $(w,\hat w) = (0,0)$. By linearity, we have for all $(v_h,\hat v_h)\in V_h^\Omega\times V_h^\Lambda$
    \begin{align*}
    \frac{1}{\tau}\int_\Omega w v_h
      +   \int_\Omega  \kappa \nabla w \cdot \nabla v_h  
      - \int_\Omega \bfU w \cdot \nabla v_h
      +\int_\Lambda\gamma \vert \partial\mathcal{D}\vert \left(\overline{w}-\hat{w}\right)\overline{v_h}
        = 0,\\
      \frac{1}{\tau} \int_\Lambda |\mathcal{D}|\hat w \hat v_h 
    + a_\Lambda(\hat w, \hat v_h)
    + b_\Lambda(\hat w, \hat v_h)
    +\int_\Lambda \gamma|\partial\mathcal{D}|\left(\hat w-\overline{w}\right)   \hat v_h
   =    0.
\end{align*}
We add the two equations and choose $(v_h,\hat v_h) = (w, \hat w)$. Using the coercivity of $a_\Lambda$ and the positivity of $b_\Lambda$, we have
\begin{align*}
    \frac{1}{\tau}\Vert w \Vert_{L^2(\Omega)}^2
    + \Vert \kappa^{1/2} \nabla w \Vert_{L^2(\Omega)}^2  
      - \int_\Omega \bfU w \cdot \nabla w
      +\Vert \gamma^{1/2} \vert \partial\mathcal{D} 
      \vert^{1/2}  (\overline{w}-\hat{w})\Vert_{L^2(\Lambda)}^2
    \\
      +\frac{1}{\tau} \Vert  |\mathcal{D}|^{1/2} \hat w \Vert_{L^2(\Lambda)}^2 
    + C_a \vert \hat w \vert_{\DG}^2
    \leq 0.
\end{align*}
In the case where $\nabla \cdot \bfU = 0$, we have
\[
\int_\Omega \bfU w \cdot \nabla w = 0,
\]
so the third term above disappears.  It is then easy to conclude that $w=\hat w = 0$ by using the coercivity of $a_\Lambda$ and the positivity of $b_\Lambda$.
%
In the case where $\bfU$ is non-divergence free, we 
write with Poincar\'e's inequality
\begin{align*}
\int_\Omega \bfU w \cdot \nabla w \leq & \Vert \bfU \Vert_{L^\infty(\Omega)} \Vert w \Vert_{L^2(\Omega)} \Vert \nabla w \Vert_{L^2(\Omega)}\\
\leq & C_0 \Vert \bfU \Vert_{L^\infty(\Omega)} \Vert \nabla w \Vert_{L^2(\Omega)}^2
\end{align*}
Therefore we have 
\[
(k_0 - C_0 \Vert \bfU \Vert_{L^\infty(\Omega)}) \Vert \nabla w \Vert_{L^2(\Omega)}^2 \leq 
\Vert \kappa^{1/2} \nabla w \Vert_{L^2(\Omega)}^2
- \int_\Omega \bfU w \cdot \nabla w.
\]
Since $\bfU$ satisfies \eqref{eq:condU},
it is then easy to conclude that $w = \hat w  = 0$. Thus, we have proved uniqueness, which is equivalent to proving existence.

Next, we show the stability bound \eqref{eq:stability}. We choose $(v_h,\hat v_h) = (c_h^n, \hat c_h^n)$ in \eqref{eq:FEMscheme}-\eqref{eq:DGscheme}. We sum the resulting equations and use coercivity of $a_\Lambda$ and positivity of $b_\Lambda$ to obtain:
\begin{align}
    \frac{1}{\tau}\int_\Omega (c_h^{n}-c_h^{n-1}) c_h^n
      + \frac{1}{\tau} \int_\Lambda |\mathcal{D}|(\hat c_h^{n}- \hat c_h^{n-1}) \hat c_h^n 
      +  (k_0 - C_0\eta_\bfU \Vert \bfU\Vert_{L^\infty(\Omega)}) \Vert \nabla c_h^n \Vert_{L^2(\Omega)}^2
      \nonumber\\
      + C_a \vert \hat c_h^n \vert_{\DG}^2 
      +\Vert \gamma^{1/2} \vert \partial\mathcal{D}\vert^{1/2} \left(\overline{c_h^{n}}-\hat{c}_h^{n}\right)\Vert_{L^2(\Lambda)}^2
%
%+\frac12 \sum_{i=1}^{N-1}  \hat U ([ \vert \mathcal{D}^{1/2}\vert \hat c_h^n]|_{s_i})^2\nonumber\\
+ \frac12 \vert \mathcal{D}(0)\vert \hat U (\hat c_h^n(0))^2
 %+\frac12 \vert \mathcal{D}(L)\vert \hat U (\hat c_h^n(L))^2
%   
\nonumber\\
        \leq  \int_\Omega f(\cdot,t^{n}) c_h^n
   + \int_\Lambda \hat f(\cdot, t^{n}) \hat c_h^n
   + |\mathcal{D}(0)|\hat U \, c^{\mathrm{in}}(t^n) \hat c_h^n(0).\label{eq:apriori1}
\end{align}
Next we observe that for the two cases we consider, we have
\begin{equation}\label{eq:twocases}
k_0 - C_0\eta_\bfU \Vert \bfU\Vert_{L^\infty(\Omega)} \geq \frac{k_0}{2}.
\end{equation}
Using Cauchy-Schwarz's inequality we have
\[
\int_\Omega f(\cdot,t^{n}) c_h^n
   + \int_\Lambda \hat f(\cdot, t^{n}) \hat c_h^n
   \leq \frac12 \Vert c_h^n\Vert_{L^2(\Omega)}^2
   + \frac12 \Vert \hat c_h \Vert_{L^2(\Lambda)}^2
   +\frac12 \Vert f(\cdot,t^n)\Vert_{L^2(\Omega)}^2
   + \frac12 \Vert \hat f(\cdot,t^n) \Vert_{L^2(\Lambda)}^2.
\]
We also have with Young's inequality:
\[
|\mathcal{D}(0)|\hat U \, c^{\mathrm{in}}(t^n)  \hat c_h^n(0)
\leq\frac14 |\mathcal{D}(0)|\hat U \,  (\hat c_h^n(0))^2
+ |\mathcal{D}(0)|\hat U \, (c^{\mathrm{in}}(t^n))^2.
\]

Therefore the bound \eqref{eq:apriori1} becomes
\begin{align}
\frac{1}{2\tau}(\Vert c_h^n\Vert_{L^2(\Omega)}^2 - \Vert c_h^{n-1}\Vert_{L^2(\Omega)}^2) 
+\frac{d_0}{2\tau}(\Vert \hat c_h^n\Vert_{L^2(\Lambda)}^2 - \Vert \hat c_h^{n-1}\Vert_{L^2(\Lambda)}^2) 
+  \frac{k_0}{2} \Vert \nabla c_h^n \Vert_{L^2(\Omega)}^2
\nonumber\\
+ C_a \vert \hat c_h^n \vert_{\DG}^2 
+\Vert \gamma^{1/2} \vert \partial\mathcal{D}\vert^{1/2} \left(\overline{c_h^{n}}-\hat{c}_h^{n}\right)\Vert_{L^2(\Lambda)}^2
%+ \frac14 \vert \mathcal{D}(0)\vert \hat U (\hat c_h^n(0))^2
\leq  
\frac12 \Vert c_h^n\Vert_{L^2(\Omega)}^2
+ \frac12 \Vert \hat c_h \Vert_{L^2(\Lambda)}^2
\nonumber\\
+ \frac12 \Vert f(\cdot,t^n)\Vert_{L^2(\Omega)}^2
+ \frac12 \Vert \hat f(\cdot,t^n) \Vert_{L^2(\Lambda)}^2
+ |\mathcal{D}(0)|\hat U \, (c^{\mathrm{in}}(t^n))^2.
\label{eq:apriori2}
\end{align}
To conclude, we multiply by $2\tau$, sum from $n=1$ to $n=m$,  use the stability of the $L^2$ projection of $\hat c^0$ in the $L^2$ norm, the approximation bound of the Scott-Zhang interpolant of $c^0$ and use Gronwall's lemma. 
% \begin{align}
%     \Vert c_h^m\Vert_{L^2(\Omega)}^2 
%     +\Vert \hat c_h^n\Vert_{L^2(\Lambda)}^2 
%       +  \tau k_0 \sum_{n=1}^m \Vert \nabla c_h^n \Vert_{L^2(\Omega)}^2
%       \nonumber\\
%       + 2 \tau C_a \sum_{n=1}^m \vert \hat c_h^n \vert_{\DG}^2 
%       +2\tau \sum_{n=1}^m \Vert \gamma^{1/2} \vert \partial\mathcal{D}\vert^{1/2} \left(\overline{c_h^{n}}-\hat{c}_h^{n}\right)\Vert_{L^2(\Lambda)}^2
%         \leq  
%         \tau \sum_{n=1}^m \Vert c_h^n\Vert_{L^2(\Omega)}^2
%    + \tau \sum_{n=1}^m \Vert \hat c_h \Vert_{L^2(\Lambda)}^2
%   \nonumber\\
% + C \Vert c^{0}\Vert_{H^s(\Omega)}^2 
% + C \Vert \hat c^{0}\Vert_{H^r(\Lambda)}^2
%        + \tau \sum_{n=1}^m \Vert f(\cdot,t^n)\Vert_{L^2(\Omega)}^2
%    + \tau \sum_{n=1}^m \Vert \hat f(\cdot,t^n) \Vert_{L^2(\Lambda)}^2
%    + 2 T |\mathcal{D}(0)|\hat U \, (c^{\mathrm{in}})^2.
% \end{align}
\end{proof}

\end{document}
